\chapter{RESULTS AND DISCUSSION}
\section{Nature of the Qualitative Findings}
The practitioner account supplies a set of specific construction-quality concerns rather than a numerical survey. The statements relate to bathroom and terrace waterproofing, PCC pavement cracks, and plaster cracking associated with masonry workmanship. Accordingly, the findings reported here concern \emph{what has been described in practice and how those concerns can be technically interpreted}. The chapter does not infer the number of affected projects, the proportion of defective work, a ranking of causes, or confirmed cost and delay impacts. The observer's identity and the associated field records remain subject to verification.

\section{W01: Waterproofing Observations and Technical Interpretation}
\subsection{Bathrooms without a continuous waterproofing layer}
The practitioner reported cases in which bathroom finishes were constructed without an adequate waterproofing layer below or behind the tiling. The associated concern is water migrating through joints or discontinuities towards adjacent rooms and finishes. The engineering issue is not simply the visual quality of the ceramic tiles. Tile and ordinary grout do not themselves establish a watertight assembly \parencite{schluter2026,tcna2023}. A wet-area system requires compatible substrates, continuous waterproofing, drains and correctly treated transitions. The visible dampness of a neighbouring room may be consistent with migration from a bathroom, but leaky plumbing, an external wall, condensation or another source must also be checked before causation is stated.

If defective waterproofing is confirmed after tiling, remedial options can involve selective tile removal, substrate assessment, correction of falls and penetrations, membrane restoration, testing and reinstatement of finishes. Each option represents rework only when tied to an actual earlier installation and corrective action; the field account does not provide quantities or records sufficient to establish that sequence for a named project.

\subsection{Workmanship, preparation and penetrations}
The practitioner associated failures with inadequate skill and supervision during membrane work. Technically, substrate cleanliness, surface soundness, primer selection, membrane thickness, continuity at corners and penetrations, and curing before covering can each influence a system's performance. Manufacturer instructions for a selected waterproofing system should govern both component compatibility and installation testing \parencite{laticrete2026}. A missing inspection record makes a failure harder to diagnose, but does not independently prove that the installation was defective.

An additional reported detail concerns holes left after removing bamboo scaffolding or temporary wall supports. Such openings may establish a potential leakage path when not soundly closed and made continuous with the surrounding wall and waterproofing. The proper repair depends on the substrate, exposure and the relevant waterproofing assembly. Merely seeing an old opening is not proof that a particular damp patch was caused by that opening; a moisture-path inspection would be needed.

\subsection{Roof terraces and drainage falls}
Insufficient drainage slope was identified as another practitioner concern. Standing water can increase the duration of moisture exposure, especially near outlets, parapets, laps and penetrations. The roof drainage path, design falls, location and condition of outlets, continuity of membrane upstands and protection from mechanical damage should be evaluated together. A roof may leak because of defective membrane detailing even when a nominal fall exists. Conversely, a small area of ponding does not automatically demonstrate membrane failure. Repairs can require restoring falls or outlets, repairing local discontinuities or reapplying a compatible membrane, followed by appropriate verification.

\subsection{Tile joints and grout}
The practitioner proposed use of tile spacers and referred to joint widths around 3--5 mm in some field installations. This is a useful description of observed practice, not a universal code threshold. Joint dimensions depend on tile size and tolerances, substrate movement, installation design and grout manufacturer's instructions. Uniform spacers may improve joint consistency and grout placement, but grout alone is not a substitute for a continuous wet-area membrane. Movement and perimeter joints also require compatible flexible detailing rather than simply filling every junction with rigid grout.

\subsection{Vertical returns, thresholds and incompatible finishes}
The field account emphasized extending waterproofing from horizontal floors up adjoining walls and at door thresholds. A return near 100 mm was mentioned as a local practice for some situations; its suitability depends on exposure, membrane system and detailed design. It must not be treated as a universal minimum, especially where shower-wall protection is required at much greater heights. The reported expression ``door seal'' should be clarified against the drawing, since it may refer to a threshold, sill, waterstop or another detail. Membrane transitions at wall--floor junctions, door thresholds and pipe penetrations demand system-specific treatment, not field coating alone \parencite{laticrete2026,schluter2026}.

The practitioner also described difficulties applying cement-based plaster directly over some sheet membranes. A smooth or otherwise incompatible membrane face may not offer a suitable key for conventional plaster; an approved render carrier, protective screed, bonding system or alternative compatible finish may be required. This is conditional on the particular material and its manufacturer specifications. No categorical conclusion about all sheet waterproofing products is justified.

\subsection{Waterproofing rework pathway}
The plausible sequence is: a waterproofing system is omitted or discontinuous; wet exposure reveals a performance concern; testing establishes the source; and a targeted correction is designed, executed and checked. Potential additional work includes removing tiles or plaster, repairing transitions or penetrations, reinstating the membrane and replacing finishes. The field material makes the concern concrete but contains no attributable labour hours, bills or completion-delay records. No numerical loss is stated.

% Safal source-only proposal; photos remain PRIVATE on Google Drive, never committed to public Git.
% Image references display placeholders unless authorized private_photos assets are mounted at build time.
\subsection{Photographic review: Lamachur residential threshold}
The researcher identifies the seven door-threshold photographs as observations at a residential house in Lamachur, Pokhara. He reports leakage following rainfall or use, exposure of an incomplete threshold detail, removal and reinstatement of previously completed work, and no observed recurrence during a later period of use. The site, observation dates, before/during/after order, engineering cause and cost records have not independently been reconciled to the photographic evidence. Hence these are researcher-reported observations, not a measured project-wide rework finding. The following plates separate visible stages without asserting an independently authenticated chronology.


\begin{figure}[p]
\centering
\begin{minipage}[t]{.48\textwidth}\centering
\IfFileExists{private_photos/SAFAL-PHOTO-W01A-01.jpg}{\includegraphics[width=\linewidth,height=5cm,keepaspectratio]{private_photos/SAFAL-PHOTO-W01A-01.jpg}}{\fbox{\parbox[c][5cm][c]{.85\linewidth}{\centering Private evidence image: SAFAL-PHOTO-W01A-01}}}
\par {\scriptsize Exposed or broken-up material next to a door threshold; SAFAL-PHOTO-W01A-01\par}\end{minipage}
\hfill
\begin{minipage}[t]{.48\textwidth}\centering
\IfFileExists{private_photos/SAFAL-PHOTO-W01A-02.jpg}{\includegraphics[width=\linewidth,height=5cm,keepaspectratio]{private_photos/SAFAL-PHOTO-W01A-02.jpg}}{\fbox{\parbox[c][5cm][c]{.85\linewidth}{\centering Private evidence image: SAFAL-PHOTO-W01A-02}}}
\par {\scriptsize Textured construction material next to wall; SAFAL-PHOTO-W01A-02\par}\end{minipage}
\par\vspace{4mm}
\begin{minipage}[t]{.48\textwidth}\centering
\IfFileExists{private_photos/SAFAL-PHOTO-W01A-03.jpg}{\includegraphics[width=\linewidth,height=5cm,keepaspectratio]{private_photos/SAFAL-PHOTO-W01A-03.jpg}}{\fbox{\parbox[c][5cm][c]{.85\linewidth}{\centering Private evidence image: SAFAL-PHOTO-W01A-03}}}
\par {\scriptsize Tiled floor and threshold with visible edge strip; SAFAL-PHOTO-W01A-03\par}\end{minipage}
\hfill
\begin{minipage}[t]{.48\textwidth}\centering
\IfFileExists{private_photos/SAFAL-PHOTO-W01A-04.jpg}{\includegraphics[width=\linewidth,height=5cm,keepaspectratio]{private_photos/SAFAL-PHOTO-W01A-04.jpg}}{\fbox{\parbox[c][5cm][c]{.85\linewidth}{\centering Private evidence image: SAFAL-PHOTO-W01A-04}}}
\par {\scriptsize Terrace-facing doorway and tiled approach; SAFAL-PHOTO-W01A-04\par}\end{minipage}
\par\vspace{4mm}
\caption{Researcher-identified Lamachur door-threshold conditions and site works (plate 1). Photographs document only visible conditions; no physical cause, case cost, or chronology is certified by image appearance alone.}\label{fig:safal-w01a-1}\end{figure}
\clearpage

\begin{figure}[p]
\centering
\begin{minipage}[t]{.48\textwidth}\centering
\IfFileExists{private_photos/SAFAL-PHOTO-W01A-05.jpg}{\includegraphics[width=\linewidth,height=5cm,keepaspectratio]{private_photos/SAFAL-PHOTO-W01A-05.jpg}}{\fbox{\parbox[c][5cm][c]{.85\linewidth}{\centering Private evidence image: SAFAL-PHOTO-W01A-05}}}
\par {\scriptsize Jamb and tiled floor with completed-looking surface work; SAFAL-PHOTO-W01A-05\par}\end{minipage}
\hfill
\begin{minipage}[t]{.48\textwidth}\centering
\IfFileExists{private_photos/SAFAL-PHOTO-W01A-06.jpg}{\includegraphics[width=\linewidth,height=5cm,keepaspectratio]{private_photos/SAFAL-PHOTO-W01A-06.jpg}}{\fbox{\parbox[c][5cm][c]{.85\linewidth}{\centering Private evidence image: SAFAL-PHOTO-W01A-06}}}
\par {\scriptsize Door threshold and tile junction after surface work; SAFAL-PHOTO-W01A-06\par}\end{minipage}
\par\vspace{4mm}
\begin{minipage}[t]{.48\textwidth}\centering
\IfFileExists{private_photos/SAFAL-PHOTO-W01A-07.jpg}{\includegraphics[width=\linewidth,height=5cm,keepaspectratio]{private_photos/SAFAL-PHOTO-W01A-07.jpg}}{\fbox{\parbox[c][5cm][c]{.85\linewidth}{\centering Private evidence image: SAFAL-PHOTO-W01A-07}}}
\par {\scriptsize Personnel and tools working around a door threshold; SAFAL-PHOTO-W01A-07\par}\end{minipage}
\par\vspace{4mm}
\caption{Researcher-identified Lamachur door-threshold conditions and site works (plate 2). Photographs document only visible conditions; no physical cause, case cost, or chronology is certified by image appearance alone.}\label{fig:safal-w01a-2}\end{figure}
\clearpage
\subsection{Other roof, parapet and finish observations}
Five further photographs supplied under the waterproofing topic show roof, parapet and finish conditions. One depicts a timber stair finish rather than a visible waterproofing system. These images may relate to different projects and do not establish a common water-entry pathway. Their project links, dates, applicable drawings and interventions require individual verification.


\begin{figure}[p]
\centering
\begin{minipage}[t]{.48\textwidth}\centering
\IfFileExists{private_photos/SAFAL-PHOTO-W01B-01.jpg}{\includegraphics[width=\linewidth,height=5cm,keepaspectratio]{private_photos/SAFAL-PHOTO-W01B-01.jpg}}{\fbox{\parbox[c][5cm][c]{.85\linewidth}{\centering Private evidence image: SAFAL-PHOTO-W01B-01}}}
\par {\scriptsize Timber stair finish showing visible surface deterioration; SAFAL-PHOTO-W01B-01\par}\end{minipage}
\hfill
\begin{minipage}[t]{.48\textwidth}\centering
\IfFileExists{private_photos/SAFAL-PHOTO-W01B-02.jpg}{\includegraphics[width=\linewidth,height=5cm,keepaspectratio]{private_photos/SAFAL-PHOTO-W01B-02.jpg}}{\fbox{\parbox[c][5cm][c]{.85\linewidth}{\centering Private evidence image: SAFAL-PHOTO-W01B-02}}}
\par {\scriptsize Parapet with downward light-coloured surface marks; SAFAL-PHOTO-W01B-02\par}\end{minipage}
\par\vspace{4mm}
\begin{minipage}[t]{.48\textwidth}\centering
\IfFileExists{private_photos/SAFAL-PHOTO-W01B-03.jpg}{\includegraphics[width=\linewidth,height=5cm,keepaspectratio]{private_photos/SAFAL-PHOTO-W01B-03.jpg}}{\fbox{\parbox[c][5cm][c]{.85\linewidth}{\centering Private evidence image: SAFAL-PHOTO-W01B-03}}}
\par {\scriptsize Roof surface discoloration at parapet junction; SAFAL-PHOTO-W01B-03\par}\end{minipage}
\hfill
\begin{minipage}[t]{.48\textwidth}\centering
\IfFileExists{private_photos/SAFAL-PHOTO-W01B-04.jpg}{\includegraphics[width=\linewidth,height=5cm,keepaspectratio]{private_photos/SAFAL-PHOTO-W01B-04.jpg}}{\fbox{\parbox[c][5cm][c]{.85\linewidth}{\centering Private evidence image: SAFAL-PHOTO-W01B-04}}}
\par {\scriptsize Sloping roof and concrete or coated edge; SAFAL-PHOTO-W01B-04\par}\end{minipage}
\par\vspace{4mm}
\caption{Roof, parapet and finish conditions recorded for separate verification (plate 1). Photographs document only visible conditions; no physical cause, case cost, or chronology is certified by image appearance alone.}\label{fig:safal-w01b-1}\end{figure}
\clearpage

\begin{figure}[p]
\centering
\begin{minipage}[t]{.48\textwidth}\centering
\IfFileExists{private_photos/SAFAL-PHOTO-W01B-05.jpg}{\includegraphics[width=\linewidth,height=5cm,keepaspectratio]{private_photos/SAFAL-PHOTO-W01B-05.jpg}}{\fbox{\parbox[c][5cm][c]{.85\linewidth}{\centering Private evidence image: SAFAL-PHOTO-W01B-05}}}
\par {\scriptsize Parapet capping with visible local joint; SAFAL-PHOTO-W01B-05\par}\end{minipage}
\par\vspace{4mm}
\caption{Roof, parapet and finish conditions recorded for separate verification (plate 2). Photographs document only visible conditions; no physical cause, case cost, or chronology is certified by image appearance alone.}\label{fig:safal-w01b-2}\end{figure}
\clearpage

\section{W02: PCC Road-Work Observations and Technical Interpretation}
\subsection{Reported cracking and slab construction}
The practitioner described cracking in PCC pavement reportedly about four to six inches thick, sometimes without distributed reinforcing bars, and drew attention to inadequate curing. These are practitioner-reported observations, not measurements extracted from authenticated designs or pavement cores. They motivate review of the approved slab thickness, anticipated vehicles and loads, support conditions, mix and curing, joint spacing and construction chronology. Because the thesis addresses RCC building projects, only such pavement within an eligible building-related work package may become an admitted research case.

\subsection{Why reinforcement cannot be treated as a universal requirement}
The absence of distributed steel does not, by itself, demonstrate faulty PCC pavement construction. FHWA describes jointed plain concrete pavement as a recognised design that need not contain distributed reinforcement, although dowels and tie bars may be included at joints to serve specific functions \parencite{fhwa2013,fhwa2009}. Of particular relevance to Nepal, Section 3202 of the Department of Roads' 2082-amended specification describes \emph{unreinforced, dowel-jointed, plain cement concrete pavement}, with slab thickness, grade and joints determined by drawings \parencite{dor2082}. There is no basis to state that any PCC slab exceeding 2.5 inches requires distributed reinforcing steel. A local slab's adequacy cannot be judged without its design loading, construction detail and project requirements.

\subsection{Cracking mechanisms to be examined}
Several explanations can produce superficially similar cracks. Early moisture loss can increase plastic shrinkage risk; restraint of drying or thermal movement can promote cracking in hardened concrete; joints placed too far apart, too late or without sufficient depth can permit uncontrolled contraction cracks. Weak or variable subgrade support, improper compaction, edge loading, unsuitable thickness for actual traffic, poor finishing or late curing may contribute to distress. In addition, construction sequencing may create weak interfaces. These hypotheses should be evaluated against crack patterns, locations, time of appearance, joint drawings, support evidence, curing records and actual loading before a cause is assigned.

The reported lack of curing deserves attention as a plausible quality concern: moisture retention is relevant to hydration, shrinkage behaviour and surface durability. Yet the narrative contains neither curing-duration records nor tests that establish a specific strength reduction. A recommendation for documented curing procedures can be made without asserting a measured deficiency at an unknown site.

\subsection{From a pavement crack to a rework event}
A crack may be monitored, sealed, repaired locally or lead to removal and replacement after engineering evaluation. Only a verified corrective intervention following earlier executed work is recorded as rework. The decision between those interventions depends on severity, structural and serviceability implications, load transfer and project acceptance requirements. The Nepal DoR specification is applicable when contractually or technically adopted for the work; it is not automatically imposed on every small internal courtyard slab. Neither the number of repairs nor their cost is determined by the supplied account.

\section{W03: Plaster Cracking Associated with Masonry Workmanship}
\subsection{Perpendicular wall junctions}
The practitioner described occasions when perpendicular brick walls met at a simple vertical interface rather than being bonded or joined by an engineered detail. Such a junction can be a location of differential movement or reduced continuity and may coincide with plaster cracking. The distinction is between a deliberately bonded intersection, an unbonded butt joint and a specified tie/connector system. Toothing is not the only legitimate construction detail, and it should not be prescribed universally. Examination of drawings, concealed connectors and crack geometry would be necessary to attribute an observed crack to poor wall-junction construction.

\subsection{Lintel and sill bands}
Omission of lintel- and sill-level bands was another practitioner concern. Where a band's presence is required by an applicable engineered drawing or code provision, its omission is a detailing compliance issue with potentially important behaviour of the wall around openings. However, the requirements vary with structural scheme and code coverage. DUDBC lists both NBC 205:2024, framed as guidance for low-rise reinforced-concrete buildings without masonry infill, and NBC 201 for buildings with masonry infill \parencite{dudbc2024}. The applicable system must be confirmed before a specific band arrangement is declared compulsory. A crack near an opening may be consistent with movement or stress concentration, but its precise cause cannot be inferred from a reported band omission alone.

\subsection{Brick moisture conditioning}
The practitioner noted that bricks are at times placed while excessively dry. Highly absorbent clay bricks can remove water from fresh mortar and impair satisfactory interaction during laying; their initial rate of absorption is therefore an important property \parencite{bailey1990,bia2026}. The appropriate pre-wetting or moisture-conditioning process is material- and weather-dependent. Saturating all brick types under all conditions is not a defensible universal recommendation. Material inspection and the approved masonry specification should instead control acceptance and preparation.

\subsection{Mortar curing and plaster cracking}
The reported inadequacy of masonry curing is also technically plausible as a workmanship concern. Cement-based mortar requires suitable moisture conditions for hydration and development of bond and strength. Early loss of water, material incompatibility, excessive plaster thickness, substrate movement, poor surface preparation and premature plaster application can each contribute to a cracking risk. A visible crack alone does not reveal which combination applies. The history of masonry execution, environmental exposure, mortar materials, plaster thickness and timing should be established before deciding whether repair is appropriate.

\subsection{Distinguishing repair from observation}
A crack photographed at the intersection of walls is a field condition. Cutting out and replacing damaged plaster, modifying an inadequately connected wall, reinstating reinforcement or another verified corrective operation may constitute rework if linked to the earlier execution. Cosmetic covering of an unresolved moving crack may recur if the substrate mechanism is not addressed. For this reason, diagnosis should precede selection of the repair method. The practitioner material does not include confirmed repair dates or quantities for any particular crack.

% Safal source-only proposal; photos remain PRIVATE on Google Drive, never committed to public Git.
% Image references display placeholders unless authorized private_photos assets are mounted at build time.
\subsection{Photographic review of brick walls and construction junctions}
Five photographs document masonry works and visible construction interfaces. The researcher reports supervising masonry and, in a distinct event, an intervention involving demolition and reconstruction. The presented images have not been linked to a dated original-work and correction sequence; therefore none is labelled a verified rework event at this stage. A marked exterior bay does not independently establish absence of concealed sill or lintel reinforcement. Drawings and inspection records govern any technical compliance claim.


\begin{figure}[p]
\centering
\begin{minipage}[t]{.48\textwidth}\centering
\IfFileExists{private_photos/SAFAL-PHOTO-W03-01.jpg}{\includegraphics[width=\linewidth,height=5cm,keepaspectratio]{private_photos/SAFAL-PHOTO-W03-01.jpg}}{\fbox{\parbox[c][5cm][c]{.85\linewidth}{\centering Private evidence image: SAFAL-PHOTO-W03-01}}}
\par {\scriptsize Unfinished masonry walls and junctions; SAFAL-PHOTO-W03-01\par}\end{minipage}
\hfill
\begin{minipage}[t]{.48\textwidth}\centering
\IfFileExists{private_photos/SAFAL-PHOTO-W03-02.jpg}{\includegraphics[width=\linewidth,height=5cm,keepaspectratio]{private_photos/SAFAL-PHOTO-W03-02.jpg}}{\fbox{\parbox[c][5cm][c]{.85\linewidth}{\centering Private evidence image: SAFAL-PHOTO-W03-02}}}
\par {\scriptsize Bricks and active site context; SAFAL-PHOTO-W03-02\par}\end{minipage}
\par\vspace{4mm}
\begin{minipage}[t]{.48\textwidth}\centering
\IfFileExists{private_photos/SAFAL-PHOTO-W03-03.jpg}{\includegraphics[width=\linewidth,height=5cm,keepaspectratio]{private_photos/SAFAL-PHOTO-W03-03.jpg}}{\fbox{\parbox[c][5cm][c]{.85\linewidth}{\centering Private evidence image: SAFAL-PHOTO-W03-03}}}
\par {\scriptsize Interior unfinished masonry wall junctions; SAFAL-PHOTO-W03-03\par}\end{minipage}
\hfill
\begin{minipage}[t]{.48\textwidth}\centering
\IfFileExists{private_photos/SAFAL-PHOTO-W03-04.jpg}{\includegraphics[width=\linewidth,height=5cm,keepaspectratio]{private_photos/SAFAL-PHOTO-W03-04.jpg}}{\fbox{\parbox[c][5cm][c]{.85\linewidth}{\centering Private evidence image: SAFAL-PHOTO-W03-04}}}
\par {\scriptsize Close-up of an RC-masonry interface, originally marked; SAFAL-PHOTO-W03-04\par}\end{minipage}
\par\vspace{4mm}
\caption{Masonry execution, junctions and building-face observations (plate 1). Photographs document only visible conditions; no physical cause, case cost, or chronology is certified by image appearance alone.}\label{fig:safal-w03-1}\end{figure}
\clearpage

\begin{figure}[p]
\centering
\begin{minipage}[t]{.48\textwidth}\centering
\IfFileExists{private_photos/SAFAL-PHOTO-W03-05.jpg}{\includegraphics[width=\linewidth,height=5cm,keepaspectratio]{private_photos/SAFAL-PHOTO-W03-05.jpg}}{\fbox{\parbox[c][5cm][c]{.85\linewidth}{\centering Private evidence image: SAFAL-PHOTO-W03-05}}}
\par {\scriptsize Masonry elevation with a marked bay; concealed band not verified; SAFAL-PHOTO-W03-05\par}\end{minipage}
\par\vspace{4mm}
\caption{Masonry execution, junctions and building-face observations (plate 2). Photographs document only visible conditions; no physical cause, case cost, or chronology is certified by image appearance alone.}\label{fig:safal-w03-2}\end{figure}
\clearpage

\section{Comparative Construction-Quality Themes}
\begin{longtable}{@{}L{0.12\textwidth}L{0.30\textwidth}L{0.46\textwidth}@{}}
\caption{Cross-category synthesis of the reported concerns}\label{tab:crosscategory}\\
\toprule Subject & Concern reported & Technical verification needed\\ \midrule\endfirsthead
\toprule Subject & Concern reported & Technical verification needed\\ \midrule\endhead
W01 & Missing or discontinuous waterproofing; inadequate falls, joints, holes or returns. & Waterproofing design and product details; substrate and membrane inspection; moisture source and performance checks.\\
W02 & PCC cracks, reported thin slabs, limited curing and questions about steel. & Drawings, actual thickness, concrete and support records, load assumptions, joint pattern, curing history and crack assessment.\\
W03 & Cracks near wall junctions or openings, dry bricks and deficient mortar curing. & Wall connection drawings, band applicability, brick absorption, mortar preparation, construction sequence and crack mapping.\\ \bottomrule
\end{longtable}

Across the three subjects, insufficient detailing, unclear inspection responsibilities, material preparation and curing can create conditions that deserve closer scrutiny. They cannot be ranked as statistically dominant causes from the supplied account. The common engineering lesson is the need to verify continuity of a designed system: a waterproofing membrane, a jointed pavement slab or a suitably connected masonry wall must be constructed and inspected as an assembly, not assessed on the superficial appearance of a finish alone \parencite{love2022,safapour2019}.

\section{Consequences and Corrective Interventions}
Potential rework in W01 may involve demolition of tiles and screed, repairs to a membrane and restoration of finishes. Potential W02 rework can include crack repair or replacement of rejected pavement sections. Potential W03 rework can include investigation, removal of plaster, correction of substrate connections where justified and reapplication. All such activities can consume extra materials, labour and time, but the account does not establish that each activity occurred at a traceable site. No frequency, direct-cost total, percentage loss or overall delay is calculated; the cost and time attribution rules in Chapter 3 indicate the records necessary to do so.

\section{Prevention and Site Quality Improvement}
For W01, prevention begins with identifying the wet-area exposure class and selecting a compatible assembly, checking substrate/falls, detailing all corners and penetrations, performing system-specified inspections or tests and preserving installation records before tiling. For W02, prevention includes verifying design loads and thickness, support preparation, joint plan and sawing sequence, proper placing/finishing and curing. For W03, prevention includes matching the masonry connection and reinforcement detail to approved drawings, controlling brick moisture and mortar use, preserving curing conditions and inspecting the substrate before applying plaster.

These measures are reasonable technical recommendations, not evidence that the current unknown sites complied or failed. They should be integrated into project inspection points and linked to documentary acceptance so any later rework claim is auditable.

\section{Discussion of Evidential Strength}
The revised qualitative discussion is substantive because the supplied narrative contains distinct practical concerns, and each has been examined against engineering mechanisms and references. It does not claim a completed local field study at a level unsupported by the available material. The unresolved observer identity, project context, dated evidence and actual repairs are significant because they separate a credible professional account from a verified case dataset. Engineering explanation is not equivalent to diagnosis, and diagnosis is not equivalent to proof of additional resources or delay. This distinction safeguards the thesis conclusions while providing a defensible direction for further field verification.
